\documentclass{ceurart}

\usepackage{lmodern}
\usepackage{mathpartir}
\usepackage{amssymb}
\usepackage{booktabs}
\usepackage{stmaryrd}
\usepackage{url}

% ── compact inference rules ──
\newcommand{\smark}[1]{\textsc{#1}}
\newcommand{\HighRisk}{\textsf{HighRisk}}
\newcommand{\LowRisk}{\textsf{LowRisk}}
\newcommand{\Denied}{\textsf{Denied}}
\newcommand{\Trust}{\textbf{Trust}}
\newcommand{\UTrust}{\textbf{UTrust}}
\newcommand{\Excess}{\textbf{Excess}}
\newcommand{\NoExcess}{\textbf{NoExcess}}
\newcommand{\wf}{\mathsf{wf}}
\newcommand{\CI}{\mathcal{P}}
\newcommand{\CItwo}{\mathcal{Q}}

\begin{document}

\copyrightyear{2026}
\copyrightclause{Copyright for this paper by its authors.
  Use permitted under Creative Commons License Attribution 4.0
  International (CC BY 4.0).}
\conference{OVERLAY 2026: 8th Workshop on Artificial Intelligence
  and Formal Verification, Logic, Automata, and Synthesis,
  co-located with AI*IA 2026}

\title{Fairness Certificates via a Lean-Backed\\Probabilistic Typed Deduction}

\author[1]{Fabio Aurelio D'Asaro}[
  orcid=0000-0002-6254-2372,
  email=fabioaurelio.dasaro@univr.it]

\author[2]{Giuseppe Primiero}[
  email=giuseppe.primiero@unimi.it]

\address[1]{Department of Human Sciences,
  University of Verona, Italy}

\address[2]{Department of Philosophy ``Piero Martinetti,''
  University of Milan, Italy}

\begin{abstract}
We present \emph{TPTND-Lean}, a Lean~4 implementation of a core fragment of
Trustworthy Probabilistic Typed Natural Deduction (TPTND).
TPTND-Lean type-checks derivation trees whose leaves are empirical observations
and whose internal nodes are introduction or elimination rules for
probabilistic connectives, producing machine-checked
\emph{fairness certificates}.
Counts, frequencies, and stored interval endpoints are exact
rational; the standard-error term used by the score-test
acceptance bands is computed by a bounded rational Newton
iteration.  We extend the original one-sample trust rules with
two-sample variants \textsc{IT2}/\textsc{IUT2} that compare two
observed populations directly, and we describe the four-stage
workflow that takes group-level summary statistics through claim
selection and derivation construction to a checked certificate.
Two case studies on public data illustrate the approach:
(i)~the ProPublica COMPAS recidivism data (6\,172 defendants,
reproducing the headline racial-bias finding), and
(ii)~HMDA mortgage-lending data for Delaware
(54\,000 originated-or-denied Black/White applications across
2022--2023, demonstrating multi-year temporal composition via
\textsc{Update} and intersectional comparisons).
Certificates are \emph{conditional}: they verify the derivation
relative to the supplied counts, the chosen statistical rule, and
the checker, not the upstream data pipeline.
The full source, data pipelines, and all derivations are
publicly available at \url{https://github.com/dasaro/tptnd-lean}.
\end{abstract}

\begin{keywords}
algorithmic fairness \sep
typed natural deduction \sep
proof checking \sep
Lean~4 \sep
trust certificates
\end{keywords}

\maketitle

% ====================================================================
\section{Introduction}\label{sec:intro}
% ====================================================================

Automated decision-making systems increasingly affect high-stakes
outcomes (criminal sentencing, credit scoring, mortgage lending),
and the question of whether their outputs are \emph{fair} across
demographic groups has become a central concern in AI
ethics~\cite{barocas2019fairness}.
Statistical tests for group-fairness properties (equalized odds,
calibration, disparate impact) are routinely applied in auditing
practice, yet the resulting claims are typically embedded in
narrative reports or Jupyter notebooks, with no formal guarantee
that the stated assumptions, data, and statistical procedures
are consistent.

Trustworthy Probabilistic Typed Natural Deduction
(TPTND)~\cite{dasaro2025tptnd} offers a proof-theoretic alternative.
TPTND is a natural-deduction calculus whose types are probability
distributions over finite output alphabets and whose terms carry
sample-size and provenance annotations.
Introduction rules for \emph{trust} and \emph{untrust} internalize
a confidence-interval membership test: if the model probability $p$
lies inside the interval $[\ell,h]$ computed from observed data, the
calculus derives $\Trust$; otherwise $\UTrust$.
Elimination rules extract the interval into the typing context,
enabling further composition.

In this paper we describe \textbf{TPTND-Lean}, a Lean~4 implementation
of a core TPTND fragment.  Our contributions are:
%
\begin{enumerate}
\item A \textbf{machine-checkable implementation} of the calculus
  (Section~\ref{sec:rules}), where every
  rule is a pure function |Derivation $\to$ CheckM Unit|
  dispatched by rule name, and all arithmetic is exact rational.
\item \textbf{Two-sample trust rules} \textsc{IT2} and \textsc{IUT2}
  (Table~\ref{tab:trust-intro}), which compare two observed
  populations without designating either as a fixed benchmark.
\item A \textbf{four-stage auditing workflow} that takes group-level
  summary statistics through claim selection and derivation
  construction to a machine-checked certificate
  (Section~\ref{sec:workflow}).
\item \textbf{Two case studies} on public data
  (Section~\ref{sec:cases}): COMPAS recidivism scoring and HMDA
  mortgage lending, producing derivation trees of depth up to~3
  that include intersectional and year-over-year comparisons.
\end{enumerate}

% ====================================================================
\section{The TPTND Calculus}\label{sec:calculus}
% ====================================================================

\smallskip\noindent
\textbf{Background.}
We treat each atomic output $a$ as an event in some implicit
finite sample space, with $P(a)$ its marginal probability.  The
output algebra mirrors standard event operations:
$\overline\alpha$ is the complement, $\alpha+\beta$ is disjoint
union (the calculus enforces $\alpha \perp_{\text{syn}} \beta$
syntactically, so $P(\alpha+\beta) = P(\alpha) + P(\beta)$),
$\alpha \times \beta$ is the joint event under an explicit
independence witness ($P(\alpha\times\beta) = P(\alpha)P(\beta)$),
and $\alpha \Rightarrow_p \beta$ encodes the conditional
$P(\beta\mid\alpha) = p$.  Two modalities of judgement accompany
distributions: \emph{expected mode} $t_n:\alpha_{\tilde p}$
records a model-side prediction that the fraction of $n$ samples
falling in $\alpha$ is $\tilde p$; \emph{frequency mode}
$t_n:\alpha_f$ records the empirical fraction $f$.  The trust
rules reconcile the two via score-test acceptance bands
$\CI(n,f,p)$.

\smallskip\noindent
We summarize the syntax, judgement forms, and key rules.  The
full calculus (31~rules across 7~families) appears
in~\cite{dasaro2025tptnd}; here we present only the rules
central to the case studies: \textsc{Obs}, \textsc{Update}, and
the trust family (Table~\ref{tab:trust-intro}).  Our main
extension is two-sample trust (\textsc{IT2}/\textsc{IUT2}).

\subsection{Syntax}\label{sec:syntax}

\emph{Outputs} are generated by
$\alpha,\beta ::= a \mid \overline\alpha \mid \alpha+\beta
  \mid \alpha\times\beta \mid \alpha\Rightarrow_p\beta$,
where $a$ ranges over a countable set of atomic names and
$p\in[0,1]$ is the conditional probability annotation on arrows.

\emph{Probability constraints} are
$c ::= p \mid [\ell,h] \mid \lnot[\ell,h]$,
with $p,\ell,h \in [0,1]$ and $\ell\le h$.
The constraint $\lnot[\ell,h]$ records that the relevant
probability lies \emph{outside} the interval; it is introduced by
\textsc{EUT}.

\emph{Terms} are $t ::= a \mid \langle t,u\rangle \mid
  \mathsf{fst}(t) \mid \mathsf{snd}(t) \mid [x]t \mid (t\cdot u)$.

A \emph{context} $\Gamma$ is a finite list of assumptions
$x : \alpha_c$ equipped with a nonempty support annotation.
Term sequents carry an explicit provenance index~$\sigma$.
For product and structural rules we write
$\Gamma\mathbin{\#_w}\Delta$ when the derivation carries an
explicit independence witness~$w$.

\subsection{Judgement forms}\label{sec:judge}

The calculus uses ten sequent forms~\cite{dasaro2025tptnd}:
\begin{center}\small
\begin{tabular}{@{}ll@{}}
$\vdash \alpha :: \mathsf{output}$ & output declaration\\
$\vdash \Gamma$ & distribution declaration\\
$\Gamma \vdash x : \alpha_c$ & identity (exact model import)\\
$\Gamma \vdash_\sigma t : \alpha$ & experiment-mode atomic claim\\
$\Gamma \vdash_\sigma t_n : \alpha_{\tilde p}$ & expected-mode claim\\
$\Gamma \vdash_\sigma t_n : \alpha_f$ & frequency-mode claim\\
$\Gamma \vdash \Trust_\CI(t_n:\alpha_f; p,[\ell,h])$
  & trust certificate\\
$\Gamma \vdash \UTrust_\CI(t_n:\alpha_f; p,[\ell,h])$
  & untrust certificate\\
$\Gamma \vdash \Excess_\CItwo(t_n:\alpha_f \succ u_m:\alpha_g; d,[\ell,h])$
  & excess certificate\\
$\Gamma \vdash \NoExcess_\CItwo(t_n:\alpha_f \succ u_m:\alpha_g; d,[\ell,h])$
  & no-excess certificate\\
\end{tabular}
\end{center}

\noindent
The one-sample acceptance band for the proportion $z$-test of a
fixed null $p$ is
\(
  \CI(n,f,p) = f \pm z_{.975}\sqrt{p(1{-}p)/n},
\)
clamped to $[0,1]$, where $z_{.975}=1.96$; the trust rules use
$p\in\CI(n,f,p)$ as the test condition.
The two-sample band, used by \textsc{IT2}/\textsc{IUT2} to test
$f{=}g$, is
\(
  \CItwo(n,m,f,g) = (f{-}g) \pm z_{.975}\sqrt{\hat p(1{-}\hat p)(1/n+1/m)},
\)
with $\hat p = (nf{+}mg)/(n{+}m)$ the pooled rate under
$H_0\!:\!p_1{=}p_2$.  We continue to write $\CI/\CItwo$ for these
bands but they are test-acceptance regions, not classical
confidence intervals.

\subsection{Global admissibility}\label{sec:admiss}

Before any rule fires, the conclusion context must satisfy $\wf(\Gamma)$:
every constraint is well-formed ($\ell\le h$),
every support is nonempty,
and for each variable the sum of exact masses and interval lower bounds
does not exceed~1.

\subsection{Running example}\label{sec:example}

Consider a simplified lending audit.
A context $\Gamma$ declares an output type $\Denied$ and
a benchmark denial rate of $20\%$: $x:\Denied_{0.2}$.
An auditor observes $n=100$ applicants from group~$A$, of whom $35$
are denied, recording this via \textsc{Obs}:
\[
\Gamma \vdash_{\{\mathit{src}_A\}} t_{100} : \Denied_{0.35}
\]
To test whether this observed rate is consistent with the $20\%$
benchmark, we apply \textsc{IUT}.
The one-sample CI is
$\CI(100, 0.35, 0.2) = 0.35 \pm 1.96\sqrt{0.2 \cdot 0.8 / 100}
  = [0.272, 0.428]$.
Since $0.2 \notin [0.272, 0.428]$, the calculus derives:
\[
\Gamma \vdash \UTrust_\CI(t_{100}:\Denied_{0.35};\; 0.2,\, [0.272, 0.428])
\]
This certificate is a self-contained artifact: anyone with the
checker can verify that the derivation is valid without re-running
the statistical test or having access to the raw data.
The next section presents the rules formally.

% ────────────────────────────────────────────────────────────────────
% Key rules (Obs, Update): inline, not separate tables
% ────────────────────────────────────────────────────────────────────

\subsection{Key rules}\label{sec:rules}

\noindent
\textbf{Observation (\textsc{Obs}).}
The fundamental data-entry rule: given a context entry supporting
term~$t$, record $n$~observations with empirical frequency~$f$:
\[
\inferrule{\sigma\neq\varnothing \quad
  \exists!\, x{:}\alpha_c\in\Gamma\text{ supporting }t \quad
  n>0 \quad nf\in\mathbb{N}}
  {\Gamma \vdash_\sigma t_n : \alpha_f}
  \;\smark{Obs}
\]
The provenance~$\sigma$ (a set of data-source identifiers) tracks
which observations contribute to a claim.

\smallskip\noindent
\textbf{Update (\textsc{Update}).}
Pools two frequency claims with disjoint provenance into a
combined weighted frequency:
\[
\inferrule{\Gamma\vdash_{\sigma_1} t_n:\alpha_f \quad
           \Gamma\vdash_{\sigma_2} t_m:\alpha_g \quad
           \sigma_1\cap\sigma_2=\varnothing}
  {\Gamma\vdash_{\sigma_1\uplus\sigma_2}
     t_{n+m}:\alpha_{\frac{nf+mg}{n+m}}}
  \;\smark{Update}
\]

\smallskip\noindent
\textbf{Identity (\textsc{Identity}).}
Imports a model probability from the context:
$\inferrule{x : \alpha_p \in \Gamma}{\Gamma \vdash x : \alpha_p}$.

\smallskip\noindent
The remaining rules (output formation, products, arrows, sum
types, structural weakening, and contraction) appear
in~\cite{dasaro2025tptnd} and are implemented in TPTND-Lean but not
exercised in the case studies below.  The Bayesian-update rules
\textsc{I-P} and \textsc{E-P} from~\cite{dasaro2025tptnd} are also
implemented and are exercised once in the COMPAS case study
(Section~\ref{sec:compas}, Table~\ref{tab:test-summary}).

% ────────────────────────────────────────────────────────────────────
% TABLE: Trust introduction (one-sample and two-sample)
% ────────────────────────────────────────────────────────────────────
\begin{table}[t]
\small
\hrule\medskip
\begin{mathpar}
\inferrule{\Gamma\vdash x:\alpha_p \quad
           \Delta\vdash_\sigma u_n:\alpha_f \\
           \CI(n,f,p)=[\ell,h] \quad p\in[\ell,h]}
  {\Gamma,\Delta\vdash\Trust_\CI(u_n:\alpha_f; p,[\ell,h])}
  \;\smark{IT}
\and
\inferrule{\Gamma\vdash x:\alpha_p \quad
           \Delta\vdash_\sigma u_n:\alpha_f \\
           \CI(n,f,p)=[\ell,h] \quad p\notin[\ell,h]}
  {\Gamma,\Delta\vdash\UTrust_\CI(u_n:\alpha_f; p,[\ell,h])}
  \;\smark{IUT}
\end{mathpar}
\medskip
\begin{mathpar}
\inferrule{\Gamma\vdash_\sigma t_n:\alpha_f \quad
           \Delta\vdash_\tau u_m:\alpha_g \\
           \sigma\cap\tau=\varnothing \\
           \CItwo(n,m,f,g)=[\ell,h] \quad 0\in[\ell,h]}
  {\Gamma,\Delta\vdash
    \Trust_\CItwo(t_n:\alpha_f; g,[\ell,h])}
  \;\smark{IT2}
\and
\inferrule{\Gamma\vdash_\sigma t_n:\alpha_f \quad
           \Delta\vdash_\tau u_m:\alpha_g \\
           \sigma\cap\tau=\varnothing \\
           \CItwo(n,m,f,g)=[\ell,h] \quad 0\notin[\ell,h]}
  {\Gamma,\Delta\vdash
    \UTrust_\CItwo(t_n:\alpha_f; g,[\ell,h])}
  \;\smark{IUT2}
\end{mathpar}
\medskip\hrule
\caption{Trust-introduction rules.
  \textsc{IT}/\textsc{IUT} are one-sample (model $p$ is a known constant).
  \textsc{IT2}/\textsc{IUT2} (new in this work) are two-sample:
  both $f$ and $g$ are estimated, and the test is whether
  $0$ lies inside the CI for $f{-}g$.}
\label{tab:trust-intro}
\end{table}

% ────────────────────────────────────────────────────────────────────
% TABLE 6: Trust elimination & comparison
% ────────────────────────────────────────────────────────────────────
\begin{table}[t]
\small
\hrule\medskip
\begin{mathpar}
\inferrule{\Gamma\vdash\Trust_\CI(u_n:\alpha_f; p,[\ell,h])}
  {\Gamma, x_u:\alpha_{[\ell,h]} \vdash_\sigma u_n:\alpha_f}
  \;\smark{ET}
\and
\inferrule{\Gamma\vdash\UTrust_\CI(u_n:\alpha_f; p,[\ell,h])}
  {\Gamma, x_u:\alpha_{\lnot[\ell,h]} \vdash_\sigma u_n:\alpha_f}
  \;\smark{EUT}
\\
\inferrule{\Gamma, x_u:\alpha_c \vdash
    \Trust_\CI(u_n:\alpha_f; p,[\ell,h]) \\
    p\in c}
  {\Gamma, x_u:\alpha_p \vdash_\sigma u_n:\alpha_{\tilde p}}
  \;\smark{ET$^{\text{ex}}$}
\end{mathpar}
\medskip
\begin{mathpar}
\inferrule{\Gamma\vdash_\sigma t_n:\alpha_f \quad
           \Delta\vdash_\tau u_m:\alpha_g \\
           \sigma\cap\tau=\varnothing \quad
           \CItwo(n,m,f,g)=[\ell,h] \quad 0\notin[\ell,h]}
  {\Gamma,\Delta\vdash\Excess_\CItwo(
    t_n:\alpha_f \succ u_m:\alpha_g;\, f{-}g,\, [\ell,h])}
  \;\smark{IEx}
\\
\inferrule{\Gamma\vdash_\sigma t_n:\alpha_f \quad
           \Delta\vdash_\tau u_m:\alpha_g \\
           \sigma\cap\tau=\varnothing \quad
           \CItwo(n,m,f,g)=[\ell,h] \quad 0\in[\ell,h]}
  {\Gamma,\Delta\vdash\NoExcess_\CItwo(
    t_n:\alpha_f \sim u_m:\alpha_g;\, f{-}g,\, [\ell,h])}
  \;\smark{INEx}
\end{mathpar}
\medskip\hrule
\caption{Trust elimination, comparison, and comparison-elimination rules.
  \textsc{EEx} and \textsc{ENEx} (not shown) re-enter the frequency
  layer under an excess or no-excess certificate, shifting the
  left observation by a model-derived interval.}
\label{tab:trust-elim}
\end{table}

% ====================================================================
\section{TPTND-Lean: Implementation in Lean 4}\label{sec:impl}
% ====================================================================

\subsection{Kernel}\label{sec:arch}

TPTND-Lean is approximately 1\,700 lines of Lean~4: a read-only kernel
(probability arithmetic, well-formedness, and one checker per
rule) surfaced through a single entry point,
$\texttt{checkDerivation}: \mathit{Derivation} \to
\mathtt{Except\,String\,Unit}$, which traverses a derivation tree
and either returns successfully, in which case the tree
\emph{is} the certificate, or yields an error message naming the
failing premise.
Three design choices reduce the trusted surface of the checker.
A probability is wrapped in a $\mathtt{Prob}$ structure carrying
proof witnesses of $0\le q\le 1$ at construction time, so an
ill-formed probability value is unrepresentable.  The checking
monad forces every rule to return either a checked result or an
explicit error, ruling out silent failures.  Rule names are a
closed inductive datatype, so the dispatcher is exhaustively
type-checked and a misspelled rule is a compile-time error.
We do not yet prove a meta-theorem connecting checker success to
an independently defined semantic validity relation: the current
trusted base consists of the Lean compiler, the
\texttt{partial}-declared dispatcher, the rational
square-root approximation (Section~\ref{sec:arith}), and the
external pipeline that produces the leaf counts.

\subsection{Rational arithmetic and provenance}\label{sec:arith}

All arithmetic is exact rational ($\mathbb{Q}$); the $[0,1]$
bounds are enforced as type-level proof obligations on every
probability value.

Square roots for the CI standard-error term are approximated by
a rational Newton iteration (15~steps), followed by rounding to
six decimal places to keep denominators bounded.
The $z$-value for the 95\% two-sided interval is the rational
$49/25 = 1.96$.
%
The square-root rounding is the one place where TPTND-Lean is not bit-exact:
the rounded value can fall on either side of the true $\sqrt{}$ by up to
$5\cdot10^{-7}$, so each CI endpoint is accurate to within
$\approx z \cdot 5\cdot10^{-7} \approx 10^{-6}$.
A derivation in which the model probability $p$ falls within
this tolerance of an interval boundary is therefore not safely
decidable by TPTND-Lean; we treat it as a known soundness boundary
rather than a guaranteed verdict, and none of the case-study
derivations sit in this regime.

\smallskip\noindent
\textbf{Provenance and probability conservation.}
The provenance index $\sigma$ has two regimes:
\textsc{I+}, \textsc{E+$_{L,R}$}, \textsc{I$\times$},
\textsc{E$\times_{L,R}$}, \textsc{E$\to$} require all three claims
to share $\sigma$ (disjoint events on one batch); \textsc{Update},
\textsc{IT2}/\textsc{IUT2}, \textsc{IEx}/\textsc{INEx} require the
two premise provenances to be \emph{disjoint} with the
conclusion's $\sigma$ their union (independent batches).
Without this discipline the textbook attack
of~\cite{dasaro2025tptnd},
\[
\dfrac{\Gamma \vdash_{\sigma_A} t_4:\alpha_{3/4} \quad
       \Gamma \vdash_{\sigma_B} t_4:\beta_{3/4}}
      {\Gamma \vdash_? t_4:(\alpha+\beta)_{3/2}}\;\smark{I+}
\]
would derive a probability of $3/2$.  In TPTND-Lean the rule layer
rejects the $\sigma_A \neq \sigma_B$ premises before any
arithmetic runs, and as a backstop the value-level primitives
\texttt{probAdd}, \texttt{probSub}, \texttt{probDiv} and
\texttt{weightedFreq} return \texttt{Option Prob} and fail
whenever a result would fall outside $[0,1]$.  A
same-$\sigma$ attack with $p+q>1$ is blocked by the value layer;
a mismatched-$\sigma$ attack with $p+q\le 1$ would slip past
the value layer and is caught only by the rule layer.

\subsection{Two-sample trust}\label{sec:it2}

The original calculus~\cite{dasaro2025tptnd} provides one-sample
rules \textsc{IT}/\textsc{IUT}, which compare an observed frequency
against a known model probability.
In practice, however, the ``model'' is often itself estimated from
data (e.g.\ the White denial rate is not a physical constant but a
sample statistic).
We add \textsc{IT2} and \textsc{IUT2}
(Table~\ref{tab:trust-intro}), which take two frequency
observations as premises and use the pooled score-test interval
$\CItwo$.
The test is: if $0\in\CItwo(n,m,f,g)$ then
\Trust{} (no significant difference); otherwise \UTrust{}.
The two-sample rules require disjoint provenances
$\sigma\cap\tau=\varnothing$, preventing a group from being
compared with (a subset of) itself.

\subsection{User workflow: producing a fairness
  certificate}\label{sec:workflow}

An auditor produces a TPTND-Lean certificate in four stages.

\smallskip\noindent
\textbf{1. Group-level summary statistics.}
The auditor extracts sample sizes and event counts for the
demographic strata of interest.  No TPTND-Lean-specific tooling is
required at this stage: SQL, pandas, or R suffice.  For COMPAS
(Section~\ref{sec:compas}) the inputs are sample sizes and
high-risk counts for each (race, sex, recidivism-status) cell;
for HMDA (Section~\ref{sec:hmda}) they are denial counts indexed
by year, race, and sex.  This is the only point at which the
auditor needs raw data.

\smallskip\noindent
\textbf{2. Claim selection.}
The fairness question fixes which TPTND-Lean rule closes the
derivation.  Common patterns:
\begin{itemize}
\item \emph{observed rate is consistent with a known benchmark
  $p$} $\to$ \textsc{IT} / \textsc{IUT};
\item \emph{two observed populations are statistically
  indistinguishable} $\to$ \textsc{IT2} / \textsc{IUT2};
\item \emph{rate increased significantly between two periods}
  $\to$ \textsc{IUT2} on year-stratified \textsc{Obs} leaves;
\item \emph{the effect persists after pooling sub-cohorts}
  $\to$ \textsc{Update} chained with one of the above.
\end{itemize}
\noindent
The English statement of the audit determines the rule signature
unambiguously, with \textsc{IT}/\textsc{IUT}/\textsc{IT2}/\textsc{IUT2}
playing the role that classical hypothesis tests play in a
statistical-software workflow.

\smallskip\noindent
\textbf{3. Derivation construction.}
The derivation tree mirrors the structure of the auditor's
English description.  Each leaf is an \textsc{Obs} node carrying
one $(n, f, \sigma)$ triple from stage~1; internal nodes apply
the rules from stage~2.  Rule arity, output-type matching, and
provenance disjointness are fixed by the rule, leaving only the
data values as free parameters.

\smallskip\noindent
\textbf{4. Verification and certification.}
A single call to \texttt{checkDerivation} traverses the tree,
applies the per-rule checker at each node, and returns either
\texttt{Ok} (the derivation \emph{is} the certificate) or an
error message that identifies the failing premise.  The
certificate is a self-contained inductive value: its validity is
independent of the data and tools used to produce it, and any
party reconstructing the derivation has the certificate.  Every
case-study derivation re-checks in well under a second.

The kernel covers stages~2--4; stage~1 remains the auditor's
responsibility.  A tabular-data front-end and a certificate
serialization format are natural next steps.

\subsection{Fairness measures expressible in TPTND}\label{sec:measures}

TPTND directly encodes any group-fairness criterion that reduces
to a difference-of-proportions test on group-conditional event
probabilities.  In the standard
taxonomy~\cite{barocas2019fairness}:

\begin{itemize}
\item \textbf{Statistical parity} ($P(\hat Y{=}1\mid A{=}a)$
  equal across groups): encoded as \textsc{IT2}/\textsc{IUT2}
  on raw acceptance rates; Tree~A in Section~\ref{sec:hmda} is
  a one-sample variant against a fixed reference.
\item \textbf{Equalized odds}~\cite{hardt2016equality}
  ($P(\hat Y{=}1\mid A{=}a, Y{=}y)$ equal across $a$ for each
  $y$): encoded as \textsc{IT2}/\textsc{IUT2} on the
  $Y$-conditioned subsets.  The ProPublica headline derivation
  (Section~\ref{sec:compas}) is the false-positive case
  ($y{=}0$); the \emph{equal-opportunity}
  relaxation~\cite{hardt2016equality} uses only $y{=}1$.
\item \textbf{Predictive parity / calibration}~\cite{chouldechova2017fair}
  ($P(Y{=}1\mid \hat Y{=}1, A{=}a)$ equal across $a$): encoded
  as \textsc{IT2} on the $\hat Y$-conditioned subsets.
  The certificates produced by TPTND-Lean should not be read as
  calibration claims: under unequal base rates and a non-perfect
  predictor, calibration / predictive parity and equalized odds
  cannot both hold~\cite{chouldechova2017fair}, so a
  COMPAS-style audit must make explicit which fairness property
  is being certified.
\end{itemize}

\noindent
Two important families fall outside this fragment.
\emph{Disparate-impact ratio} (an $80\%$ rule on
$P(\hat Y{=}1\mid A{=}0)/P(\hat Y{=}1\mid A{=}1)$) is
multiplicative and would require an Excess-style family with
division.  \emph{Individual} and \emph{counterfactual}
fairness require causal structure, addressed by the
causal-label extensions of
TPTND~\cite{ceragioli2025causal1,ceragioli2025causal2}.

% ====================================================================
\section{Case Studies}\label{sec:cases}
% ====================================================================

\smallskip\noindent
\textbf{What the certificates verify.}
The certificates produced below are \emph{conditional}.  The
checker verifies that the supplied counts and the chosen rule
form a valid TPTND derivation; it does not certify that the
upstream data extraction (filtering, demographic encoding,
outcome construction, hypothesis selection) was done correctly.
Those steps remain part of the external audit trail.  We also
inherit the standard caveats of the underlying statistical
tests (binomial sampling, independence, large-sample
approximation, no multiple-testing correction).

\subsection{COMPAS: recidivism risk scoring}\label{sec:compas}

The first case study uses the ProPublica COMPAS
dataset~\cite{angwin2016machine}
(\texttt{compas-scores-two-years.csv}, filtered to 6\,172
defendants by the published criteria).
We construct 12~Lean-validated derivation trees, including:

\smallskip\noindent
\textbf{ProPublica's headline (\UTrust{}).}\;
ProPublica's article reports false-positive rates of $44.9\%$ for
Black and $23.5\%$ for White non-recidivists~\cite{angwin2016machine}.
On the published filter (\texttt{days\_b\_screening\_arrest}~$\in[-30,30]$,
\texttt{is\_recid}~$\neq -1$, \texttt{c\_charge\_degree}~$\neq$~\texttt{O},
\texttt{score\_text}~$\neq$~\texttt{N/A}, yielding 6\,172 cases) the
corresponding observed rates are $641/1514 \approx 42.3\%$ for Black
and $282/1281 = 94/427 \approx 22.0\%$ for White, preserving the
${\sim}1.9\times$ disparity.
The White FPR is itself a sample statistic, so the statistically
more appropriate comparison uses \textsc{IUT2}; we re-derive the
two-sample variant in the HMDA case study below
(Section~\ref{sec:hmda}, Trees~B and~C) and present the
\textsc{IUT} version here only for continuity with ProPublica's
fixed-benchmark framing.
Let $\Gamma_{BN}$ carry the non-exact assumption $\HighRisk_{[0,1]}$
for the Black cohort and $\Gamma^{\text{mdl}}_{WN} \vdash x_{WN}:\HighRisk_{94/427}$
import the White benchmark; write
$\Theta_{FP} = \Gamma^{\text{mdl}}_{WN}, \Gamma_{BN}$ and
$\sigma_B = \sigma_m \uplus \sigma_f$.
The TPTND-Lean-verified derivation is:
%
\[
\scriptstyle
\inferrule{
  \Gamma_{BN}\vdash_{\sigma_m} u_{1168}:\HighRisk_{255/584} \\
  \Gamma_{BN}\vdash_{\sigma_f} u_{346}:\HighRisk_{131/346}
}{
  \Gamma_{BN}\vdash_{\sigma_B} u_{1514}:\HighRisk_{641/1514}
}\;\smark{Update}
\]
\[
\scriptstyle
\inferrule{
  \Gamma^{\text{mdl}}_{WN}\vdash x_{WN}:\HighRisk_{94/427} \\
  \Gamma_{BN}\vdash_{\sigma_B} u_{1514}:\HighRisk_{641/1514} \\
  94/427 \notin \CI(1514,\, 641/1514,\, 94/427) = [0.4025, 0.4443]
}{
  \Theta_{FP}\vdash \UTrust_\CI(u_{1514}:\HighRisk_{641/1514};\;
    94/427,\, [0.4025, 0.4443])
}\;\smark{IUT}
\]
\[
\scriptstyle
\inferrule{
  \Theta_{FP}\vdash \UTrust_\CI(u_{1514}:\HighRisk_{641/1514};\;
    94/427,\, [0.4025, 0.4443])
}{
  \Theta_{FP}, x_{BN}:\HighRisk_{\lnot[0.4025, 0.4443]}
   \vdash_{\sigma_B} u_{1514}:\HighRisk_{641/1514}
}\;\smark{EUT}
\]
%
The resulting term claim, with $\HighRisk_{\lnot[0.4025, 0.4443]}$
recorded in its typing context, is what a regulator would receive
and re-check.

\smallskip\noindent
\textbf{Bayesian belief update (depth~2).}\;
A pilot auditor entertaining three hypotheses about the FPR
($a_1{=}1/6$, $a_2{=}1/3$, $a_3{=}1/2$) with prior weights
$b_i{=}a_i$ (summing to~1) observes $s{=}2$ flagged among
$n{=}5$ pilot defendants.  \textsc{I-P} introduces the prior
family from three exact identity premises; \textsc{E-P}
produces the posterior under hypothesis~$H_j$:
\[
\textstyle
P(H_j\mid d) =
\frac{a_j^s (1-a_j)^{n-s}\, b_j}
     {\sum_i a_i^s (1-a_i)^{n-s}\, b_i}.
\]
For~$H_3$ this gives $729/1366 \approx 53.4\%$ (up from $50\%$)
while $H_1$ drops to $125/1366 \approx 9.2\%$.  The computation
is exact rational and the conclusion is a verified posterior
usable as a model premise in a downstream \textsc{IT}.

\smallskip\noindent
\textbf{A genuine non-finding (\Trust{}).}\;
Among Black recidivists, the female lower-risk rate is $62/203$
and the male benchmark is $411/1458 = 137/486$.
Since $\CI(203, 62/203, 137/486) = [0.2435, 0.3673]$
and $137/486 \approx 0.2819 \in [0.2435, 0.3673]$,
the calculus derives \Trust{} and \textsc{ET$^{\text{ex}}$}
re-enters the expected layer with the trusted exact value.

\subsection{HMDA: mortgage lending fairness}\label{sec:hmda}

The second case study uses Home Mortgage Disclosure Act (HMDA)
data for Delaware, years 2022 and 2023, downloaded from the
CFPB Data Browser.\footnote{%
\url{https://ffiec.cfpb.gov/data-browser/?category=states&items=DE}}
HMDA is a federally mandated dataset under Regulation~C
(12~CFR~1003), making the ``transferable certificate'' framing
directly grounded in actual regulatory practice.

After restricting to originated ($\texttt{action\_taken}=1$) and
denied ($\texttt{action\_taken}=3$) applications, we obtain
the summary in Table~\ref{tab:hmda-data}.

\begin{table}[t]
\small
\centering
\begin{tabular}{@{}llrrr@{}}
\toprule
Year & Group & Denied & Total & Rate \\
\midrule
2022 & White & 4\,914  & 24\,614 & 19.96\% \\
     & Black & 2\,409  & 7\,142  & 33.73\% \\
2023 & White & 3\,908  & 17\,355 & 22.52\% \\
     & Black & 2\,035  & 5\,394  & 37.73\% \\
\midrule
2023 & Black Male   &  741 & 1\,825 & 40.60\% \\
     & White Male   & 1\,466 & 5\,450 & 26.90\% \\
     & Black Female &  909 & 2\,341 & 38.83\% \\
     & White Female & 1\,234 & 4\,634 & 26.63\% \\
\bottomrule
\end{tabular}
\caption{HMDA Delaware denial rates, 2022--2023.}
\label{tab:hmda-data}
\end{table}

We build six derivation trees, organised in four cases that
each highlight a different aspect of the certificate framework:

\smallskip\noindent
\textbf{Tree~A: reference-rate trust (depth~2).}\;
An \textsc{Identity} leaf pins a $20\%$ reference denial rate; an
\textsc{Obs} leaf records the White DE~2022 frequency
($4914/24614$).
\textsc{IT} derives \Trust{}: the observed rate is consistent
with the reference.
Every assumption (reference value, data source, sample size,
statistical test) is explicit in the certificate and
machine-checkable.

\smallskip\noindent
\textbf{Tree~B: temporal composition (depth~3).}\;
Two \textsc{Obs} leaves per racial group (one for 2022, one for
2023) are combined via \textsc{Update} into pooled frequencies.
The derivation structure is:
\[
\scriptstyle
\inferrule{
  \inferrule{
    \Gamma\vdash_{\sigma_{22}} t:\Denied_{2409/7142} \quad
    \Gamma\vdash_{\sigma_{23}} t:\Denied_{2035/5394}
  }{
    \Gamma\vdash_{\sigma_B} t_{12536}:\Denied_{4444/12536}
  }\;\smark{Upd.}
  \\\\
  \inferrule{
    \Gamma\vdash_{\tau_{22}} u:\Denied_{4914/24614} \quad
    \Gamma\vdash_{\tau_{23}} u:\Denied_{3908/17355}
  }{
    \Gamma\vdash_{\tau_W} u_{41969}:\Denied_{8822/41969}
  }\;\smark{Upd.}
}{
  \Gamma\vdash\UTrust_\CItwo(\ldots;\, [0.1357, 0.1528])
}\;\smark{IUT2}
\]
Since $0 \notin [0.1357, 0.1528]$, it yields \UTrust{}.  The
$\UTrust_\CItwo$ certificate is itself transferable: a regulator
with the kernel can re-check the entire depth-3 chain
(\textsc{Obs}$\to$\textsc{Update}$\to$\textsc{IUT2}) without
access to the raw data.

\smallskip\noindent
\textbf{Trees~C.1/C.2: intersectional certificates (depth~2).}\;
Two auditors independently certify Black-vs-White denial gaps for
males ($741/1825$ vs.\ $1466/5450$, $\CItwo=[0.1126,0.1614]$) and
females ($909/2341$ vs.\ $1234/4634$, $\CItwo=[0.0990,0.1449]$);
both derive \UTrust{}.  A regulator verifies both certificates
with the same kernel without needing access to either auditor's
raw data.

\smallskip\noindent
\textbf{Trees~D.1/D.2: year-over-year monitoring (depth~2).}\;
Black denial rose 2022$\to$2023 from $33.7\%$ to $37.7\%$
($\CItwo=[0.0230,0.0568]$, \UTrust{}); White from $20.0\%$ to
$22.5\%$ ($\CItwo=[0.0176,0.0334]$, \UTrust{}); both increases
are statistically significant.  Persistent national-level racial
disparities in mortgage lending are documented by the Urban
Institute~\cite{urban2023}.  Dated certificates accumulate in a
regulatory audit trail.

\subsection{Summary}

Table~\ref{tab:test-summary} summarizes all Lean-validated
derivation trees across both case studies.

\begin{table}[t]
\small\centering
\begin{tabular}{@{}lllrc@{}}
\toprule
Dataset & Derivation & Rules & Depth & Result \\
\midrule
COMPAS & Caucasian FNR & \smark{IT} & 2 & \Trust \\
       & UPDATE chain & \smark{Update} & 2 & pooled \\
       & AA FNR \S4.3 & \smark{IT} & 2 & \Trust \\
       & AA FPR gender & \smark{IUT} & 2 & \UTrust \\
       & ProPublica headline & \smark{Upd.+IUT} & 3 & \UTrust \\
       & Bayesian & \smark{I-P+E-P} & 2 & posterior \\
       & Racial FPR & \smark{IEx} & 2 & Excess \\
       & AA FNR gender & \smark{INEx} & 2 & NoExcess \\
       & \multicolumn{2}{l}{+ 4 derivations with \smark{IUT2, IT2}} & 2 & varied \\
\midrule
HMDA   & Reference-rate & \smark{IT} & 2 & \Trust \\
       & Temporal pool & \smark{Upd.+IUT2} & 3 & \UTrust \\
       & Intersect.\ (M) & \smark{IUT2} & 2 & \UTrust \\
       & Intersect.\ (F) & \smark{IUT2} & 2 & \UTrust \\
       & Y-o-Y Black & \smark{IUT2} & 2 & \UTrust \\
       & Y-o-Y White & \smark{IUT2} & 2 & \UTrust \\
\bottomrule
\end{tabular}
\caption{Lean-validated derivation trees (all pass
  \texttt{checkDerivation}).}
\label{tab:test-summary}
\end{table}


% ====================================================================
\section{Related Work}\label{sec:related}
% ====================================================================

TPTND was introduced by D'Asaro, Genco, and
Primiero~\cite{dasaro2025tptnd}; a companion $\lambda$-calculus
formulation appears in~\cite{genco2023lambda}, and extensions
with causal labels by Ceragioli and
Primiero~\cite{ceragioli2025causal1,ceragioli2025causal2}.
The \textsc{Brio} tool~\cite{brio2023} evaluates bias via distance
metrics but does not produce derivation trees.
The certifying-algorithms pattern of separating producer from
verifier is classical~\cite{necula1997pcc,mcconnell2011certifying}.
VeriFair~\cite{bastani2019verifair} and
FairSquare~\cite{albarghouthi2017fairsquare} verify fairness
properties \emph{a priori} (of the model); TPTND-Lean instead produces
\emph{a posteriori} certificates from observed data, checkable
without model access.  Lean~4~\cite{moura2021lean4} has been used
for large-scale mathematical formalization but, to our knowledge,
not previously for algorithmic-fairness auditing.  Statistical
toolkits such as Fairlearn~\cite{bird2020fairlearn},
AIF360~\cite{bellamy2019aif360}, and
What-If~\cite{wexler2019whatif} compute metrics but do not
produce formally verified certificates: TPTND-Lean's architectural
difference is the separation of the (potentially biased)
\emph{producer} from the (small, auditable) \emph{checker}.

% ====================================================================
\section{Conclusion}\label{sec:concl}
% ====================================================================

We have presented TPTND-Lean, a Lean~4 implementation of a core TPTND
fragment that type-checks derivation trees over real-world
fairness data.
The implementation adds two-sample trust rules (\textsc{IT2}/\textsc{IUT2}),
specifies a four-stage workflow that takes raw group-level
statistics through claim selection and derivation construction
to a verified certificate, and exhibits the workflow on two
public datasets spanning criminal justice and mortgage lending.

The main limitation is that the \texttt{Prob} type is clamped
to $[0,1]$, which loses information when the two-sample CI
for a negative difference falls entirely below zero.
In practice this is handled by ordering the premises so
that the larger rate appears first.
Future work includes signed-difference probability handling,
connecting to the causal-label extensions
of~\cite{ceragioli2025causal1}, and a front-end that constructs
derivation trees from tabular data so domain experts can produce
certificates without writing Lean.

% ====================================================================
% References
% ====================================================================
\begin{thebibliography}{99}

\bibitem{dasaro2025tptnd}
F.\,A.\ D'Asaro, F.\,A.\ Genco, G.\ Primiero,
``Checking trustworthiness of probabilistic computations in a
typed natural deduction system,''
\emph{J.\ Logic and Computation}, vol.~35, no.~6, 2025.
\url{https://doi.org/10.1093/logcom/exaf003}.

\bibitem{genco2023lambda}
F.\,A.\ Genco, G.\ Primiero,
``A typed lambda-calculus for establishing trust in probabilistic programs,''
arXiv:2302.00958, 2023.

\bibitem{ceragioli2025causal1}
L.\ Ceragioli, G.\ Primiero,
``A proof system with causal labels (Part~I):
checking individual fairness and intersectionality,''
in \emph{OVERLAY 2025}, CEUR-WS vol.~4142, 2025, pp.~173--179.
\url{https://ceur-ws.org/Vol-4142/paper19.pdf}.

\bibitem{ceragioli2025causal2}
L.\ Ceragioli, G.\ Primiero,
``A proof system with causal labels (Part~II):
checking counterfactual fairness,''
in \emph{OVERLAY 2025}, CEUR-WS vol.~4142, 2025, pp.~181--187.
\url{https://ceur-ws.org/Vol-4142/paper20.pdf}.

\bibitem{brio2023}
G.\ Coraglia, F.\,A.\ D'Asaro, F.\,A.\ Genco \emph{et al.},
``BRIOxAlkemy: a bias detecting tool,''
in \emph{BEWARE@AI*IA 2023}, CEUR-WS vol.~3615, 2023.

\bibitem{angwin2016machine}
J.\ Angwin, J.\ Larson, S.\ Mattu, L.\ Kirchner,
``Machine bias,''
\emph{ProPublica}, May 23, 2016.

\bibitem{urban2023}
L.\ Goodman, J.\ Zhu,
``What different denial rates can tell us about racial
disparities in the mortgage market,''
\emph{Urban Institute}, 2023.

\bibitem{barocas2019fairness}
S.\ Barocas, M.\ Hardt, A.\ Narayanan,
\emph{Fairness and Machine Learning: Limitations and Opportunities},
fairmlbook.org, 2019.

\bibitem{hardt2016equality}
M.\ Hardt, E.\ Price, N.\ Srebro,
``Equality of opportunity in supervised learning,''
in \emph{NeurIPS}, 2016, pp.~3323--3331.

\bibitem{chouldechova2017fair}
A.\ Chouldechova,
``Fair prediction with disparate impact: a study of bias in
recidivism prediction instruments,''
\emph{Big Data}, vol.~5, no.~2, pp.~153--163, 2017.

\bibitem{necula1997pcc}
G.\,C.\ Necula,
``Proof-carrying code,''
in \emph{POPL '97}, ACM, 1997, pp.~106--119.

\bibitem{mcconnell2011certifying}
R.\,M.\ McConnell, K.\ Mehlhorn, S.\ N\"aher, P.\ Schweitzer,
``Certifying algorithms,''
\emph{Computer Science Review}, vol.~5, no.~2, pp.~119--161, 2011.

\bibitem{bastani2019verifair}
O.\ Bastani, X.\ Zhang, A.\ Solar-Lezama,
``Probabilistic verification of fairness properties via concentration,''
in \emph{Proc.\ ACM Program.\ Lang.\ (OOPSLA)}, 2019,
article~118.

\bibitem{albarghouthi2017fairsquare}
A.\ Albarghouthi, L.\ D'Antoni, S.\ Drews, A.\,V.\ Nori,
``FairSquare: probabilistic verification of program fairness,''
in \emph{Proc.\ ACM Program.\ Lang.\ (OOPSLA)}, 2017,
article~80.

\bibitem{moura2021lean4}
L.\ de Moura, S.\ Ullrich,
``The Lean~4 theorem prover and programming language,''
in \emph{CADE-28}, LNCS vol.~12699, 2021, pp.~625--635.

\bibitem{bird2020fairlearn}
S.\ Bird \emph{et al.},
``Fairlearn: a toolkit for assessing and improving fairness in AI,''
Microsoft Research, Tech.\ Rep.\ MSR-TR-2020-32, 2020.

\bibitem{bellamy2019aif360}
R.\,K.\,E.\ Bellamy \emph{et al.},
``AI Fairness 360: an extensible toolkit for detecting and mitigating
algorithmic bias,''
\emph{IBM J.\ Res.\ Dev.}, vol.~63, no.~4/5, 2019.

\bibitem{wexler2019whatif}
J.\ Wexler \emph{et al.},
``The What-If Tool: interactive probing of machine learning models,''
\emph{IEEE TVCG}, vol.~26, no.~1, pp.~56--65, 2020.

\end{thebibliography}

\end{document}
